% Compiles on its own (pulls ../main.tex's preamble via subfiles) or as part of the book.
\documentclass[../main.tex]{subfiles}
\begin{document}

\chapter{Training: The Loss and SimOTA}

The detector runs but cannot learn. To learn it needs a \emph{loss} --- one number measuring how
wrong its predictions are --- which we minimise by gradient descent. This chapter is the theory
gate for that machinery. It is also the one place we deliberately \emph{borrow} rather than build,
and the chapter explains both the mechanism and that choice.

\section{The problem the loss must solve}

Recall the head makes 6720 predictions per image, one per grid cell, each with a box, an
objectness score, and a class score. A training image has maybe one or two real pedestrians. So
before we can score anything, we face a question with no obvious answer:

\begin{quote}
Of the 6720 cells, \emph{which ones} are supposed to predict the pedestrian?
\end{quote}

This is \textbf{label assignment}, and it is the heart of modern detector training. Get it wrong
and the loss is meaningless. The naive answers are both bad: ``only the single cell at the
pedestrian's centre'' wastes all the other nearby cells that could also learn the object; ``every
cell overlapping the box'' floods training with near-duplicate positives. We need something
adaptive.

\section{SimOTA, in words}

YOLOX's answer is \textbf{SimOTA} (Simplified Optimal Transport Assignment). For one image it works
like this:

\begin{enumerate}
  \item \textbf{Center prior.} For each ground-truth pedestrian, gather the cells whose centre
        falls inside the box (or within a small radius of the box centre). Only these are
        candidates --- a cell far away should never be responsible for this object.
  \item \textbf{Cost matrix.} For every (candidate cell, pedestrian) pair, compute a cost:
        \[
          \text{cost} = \text{classification cost} + \lambda \cdot \text{box cost},
        \]
        where the box cost is $-\log(\text{IoU})$ between the cell's predicted box and the
        pedestrian, and the classification cost is how wrong the cell's class score is. Low cost
        = this cell is already good at predicting this pedestrian.
  \item \textbf{Dynamic $k$.} Here is the clever part. For each pedestrian, look at its top few
        candidate IoUs and \emph{sum} them to get an integer $k$ --- an easy, well-covered object
        gets a larger $k$ (assign it more cells), a hard or tiny one gets a smaller $k$. Then pick
        the $k$ lowest-cost cells as that pedestrian's positives. The number of positives per
        object is decided \emph{per object, from the data}, not fixed by hand.
  \item \textbf{Resolve conflicts.} If a cell got assigned to two pedestrians, keep the
        lower-cost one. Everything not assigned is background.
\end{enumerate}

The output is a \emph{foreground mask} (which cells are positives) and, for each positive, which
pedestrian it must predict. Now the loss is well-defined.

\section{The three-part loss}

With assignment done, the loss has three components, summed:

\begin{itemize}
  \item \textbf{Box loss} (positives only): $1-\text{IoU}$ between each positive cell's predicted
        box and its assigned pedestrian. Drives the boxes to fit.
  \item \textbf{Objectness loss} (all cells): binary cross-entropy pushing positives toward
        ``object here'' and everything else toward ``background.'' This is what makes the sea of
        random boxes recede at inference.
  \item \textbf{Classification loss} (positives only): binary cross-entropy on the class scores of
        the positive cells. (With one class this is small, but it stays for generality.)
\end{itemize}

The total is normalised by the number of positives so images with more pedestrians do not dominate.

\section{Why we borrow this, and what that protects}

We port SimOTA and the loss from the official YOLOX rather than reimplement them. The reasoning is
methodological. This machinery is \emph{held identical across every experiment in the study} --- it
is not a variable we ablate. A subtle bug in the cost matrix or the dynamic-$k$ logic would shift
which cells become positives, and therefore move every result in the table \emph{by the same
hidden amount}, in a way no single experiment could reveal. Reimplementing it would risk exactly
that, for no scientific gain, and would cost days of debugging. So we borrow the fixed part and
spend our care on the parts we vary --- the fusion.

\begin{lesson}
There is a live interaction to remember (it recurs in later chapters). Label assignment depends on
the predictions, and fusion changes the predictions. So changing the fusion can change which cells
SimOTA marks as positive --- producing an accuracy difference that is really an assignment
difference, not a feature-quality difference. When a fusion ablation moves unexpectedly, inspect
the number and location of positives before telling an architectural story about it.
\end{lesson}

\section{How we verify borrowed code}

Borrowed is not the same as trusted. Two checks give us confidence the port is wired correctly:
\begin{enumerate}
  \item the loss is a finite scalar and gradients flow to every parameter; and
  \item on a single image trained in a tight loop, the loss \emph{decreases} and the model
        \emph{overfits} it --- if a model cannot memorise one image, its training is broken. This
        ``overfit one batch'' test is the cheapest, most reliable smoke test in deep learning, and
        it is the exit check for this chapter.
\end{enumerate}
\end{document}
